% Generated draft: requires mathematical and source review.
\documentclass[11pt]{article}
\usepackage[margin=25mm]{geometry}
\usepackage{fontspec}
\setmainfont{Latin Modern Roman}
\usepackage{amsmath,amssymb,mathtools,unicode-math}
\setmathfont{Latin Modern Math}
\usepackage{enumitem,needspace,xurl,hyperref}
\setlength{\parindent}{0pt}
\setlength{\parskip}{5pt}
\setlength{\emergencystretch}{3em}
\newcommand{\N}{\mathbb{N}}
\newcommand{\Z}{\mathbb{Z}}
\newcommand{\Q}{\mathbb{Q}}
\newcommand{\R}{\mathbb{R}}
\newcommand{\C}{\mathbb{C}}
\newcommand{\F}{\mathbb{F}}
\newcommand{\A}{\mathbb{A}}
\newcommand{\PP}{\mathbb{P}}
\newcommand{\E}{\mathbb{E}}
\newcommand{\eps}{\varepsilon}
\newcommand{\dd}{\,\mathrm{d}}
\begin{document}
\section*{UnsolvedMath ID: 30001037}
\textbf{Source problem code:} catalogue entry
\subsection{Definitions and mathematical statement}
Under the usual orientation-preserving convention, an interval exchange transformation is a bijection \(T:[0,1)\to[0,1)\) for which there is a finite partition into half-open intervals \(I_1,\ldots,I_r\) such that \(T|_{I_i}\) is a translation and the image intervals partition \([0,1)\). These transformations form a group \(G\) under composition. This definition fixes a conventional reading; the extraction itself does not specify whether transformations with flips, endpoint identifications, or other variants are included.

The source asks two broad questions: describe the group-theoretic structure of \(G\), and characterize the abstract groups \(H\) for which there exists an injective homomorphism \(H\to G\). No finite-generation, countability or other restriction on \(H\) is stated. The structure question does not name a particular invariant or proposed property, and the subgroup question does not supply a conjectured classification. Accordingly these are retained as research directions, not converted into an invented universal assertion about amenability, free subgroups, or another property.
\subsection{Short English statement}
Study the algebraic structure of the interval-exchange group and determine which abstract groups embed in it. The conventional no-flips definition is explained, while the source's broad classification scope and missing variant conventions remain explicit.
\subsection{Sources}
\begin{itemize}
\item Local UnsolvedMath registry, permanent ID \(30001037\), source code \texttt{catalogue entry}, ``Subgroup Structure of Interval Exchange Transformations.'' Registry curation: Ulam AI / UnsolvedMath Contributors, CC BY 4.0; underlying sources retain their supplied provenance.
\item \emph{Geometric Group Theory, Hyperbolic Dynamics and Symplectic Geometry} (2009), Oberwolfach Reports, pp.\ 2337--2406, DOI \url{https://doi.org/10.4171/owr/2008/42}.
\item Supplied ``Dahmani--Fujiwara--Guirardel'': \url{https://arxiv.org/abs/1101.5909}.
\item Supplied ``Dahmani, survey'': \url{https://numdam.org/articles/10.5802/wbln.27/}.
\item Supplied literature triage, \texttt{literature\_checked\_at}: 2026-08-21: labeled open; says primary and survey literature provides substantial subgroup obstructions but no full structural or subgroup classification resolving these questions. This is supplied dated curation, not independently verified.
\end{itemize}
\end{document}
